%% kapa: exam=1 round=31 year=2020 period=2 form=A subject=accounting subject_ko=회계학 seq=30 number=70
%% source: Q-Net 감정평가사 기출문제 (cst00309) artlSeq=5207219 file=extracted/1008/5207219/A형/A형/2020년도 제31회 감정평가사 1차 시험_2교시_A형.pdf
%% conversion: kordoc 4.24.51 -> tools/gen1.py (auto); review=proofread-sonnet; flags=none
\begin{question}{70}
20x1년 1월 1일 (주)감평은 (주)한국이 동 일자에 발행한 사채(액면금액 ￦1,000,000, 액면이자율 연 4\%, 이자는 매년 말 지급)를 ￦896,884에 취득하였다. 취득 당시 유효이자율은 연 8\%이다. 20x1년 말 동 사채의 이자수취 후 공정가치는 ￦925,000이며, 20x2년 초 ￦940,000에 처분하였다. (주)감평의 동 사채 관련 회계처리에 관한 설명으로 옳지 않은 것은? (단, 계산금액은 소수점 첫째자리에서 반올림하며, 단수차이로 인한 오차가 있으면 가장 근사치를 선택한다.)
\begin{choices}
  \item 당기손익-공정가치(FVPL) 측정 금융자산으로 분류하였을 경우, 20x1년 당기순이익은 ￦68,116 증가한다.
  \item 상각후원가(AC) 측정 금융자산으로 분류하였을 경우, 20x1년 당기순이익은 ￦71,751 증가한다.
  \item 기타포괄손익-공정가치(FVOCI) 측정 금융자산으로 분류하였을 경우, 20x1년 당기순이익은 ￦71,751 증가한다.
  \item 상각후원가(AC) 측정 금융자산으로 분류하였을 경우, 20x2년 당기순이익은 ￦11,365 증가한다.
  \item 기타포괄손익-공정가치(FVOCI) 측정 금융자산으로 분류하였을 경우, 20x2년 당기순이익은 ￦15,000 증가한다.
\end{choices}
\end{question}
\answer{5}
